\subsection{Basic Set Operations}

Now that we have a basic understanding of sets, we will now learn how to work with multiple sets. What do we do if we want to compare elements across different sets? For instance, how can we ``add" and ``subtract" elements of different sets, or find elements that they have in common? This subsection aims to introduce the fundamental operations on sets that help us answer these questions.

\begin{defbox}
   Let $A$ and $B$ be two sets. Define:
   \begin{itemize}
      \item Their \textbf{union}, denoted $A \cup B$, as the set $\{x \colon x \in A \text{ or } x \in B\}$
      \item Their \textbf{intersection}, denoted $A \cap B$, as the set $\{x \colon x \in A \text{ and } x \in B\}$
      \item Their \textbf{difference}, denoted $A \setminus B$, as the set $\{x \colon x \in A \text{ and } x \notin B\}$.
   \end{itemize}
\end{defbox}

If $A$ is an ambient set (i.e. there are other sets that are contained in $A$), then if $B \subseteq A$ we say the \textbf{complement} of $B$ relative to $A$, denoted $B^c$, is the difference:
\[ B^c := A \setminus B \]

\begin{example}
   Within $\ZZ$, define $A := \{n \in \ZZ \colon n \text{ is even}\}$. We have that the complement of $A$ relative to $\ZZ$, $A^c := \{n \in \ZZ \colon n \text{ is odd}\}$. 
\end{example}

\begin{defbox}
   For two elements $a$ and $b$ ($a$ and $b$ are not necessarily sets), we define their \bold{ordered pair}, denoted $(a,b)$, to be the tuple satisfying the property:
   \[ (a,b)=(c,d) \quad\text{if and only if}\quad a=c \text{ and } b=d \]
In particular, if \(a\neq b\), then \((a,b)\neq(b,a)\).$(a,b) \neq (b,a)$ whenever $a \neq b$. 
\end{defbox}

We will use this notion of ordered pairs to help us with the following important definition:

\begin{defbox}
   Let $A$, $B$ be two sets. We define their \textbf{Cartesian product}, denoted $A \times B$, as the set of ordered pairs $\{(a,b) \colon a \in A \text{ and } b \in B\}$.
\end{defbox}

A notable Cartesian product is the Euclidean plane $\RR^2 := \RR \times \RR$. You have already seen this plane defined as the $xy$-plane, where each ordered pair $(x,y)$ represents the $x$- and $y$-coordinates of a point on the plane.

\begin{example}
   Let $A := \{0, 1, 2\}$ and $B := \{a, b\}$. Then their Cartesian product $A \times B$ is the set $\{(0,a), (0,b), (1,a), (1,b), (2,a), (2,b)\}$.
\end{example}

\begin{remark}
   In general, if $A$ has $n$ elements and $B$ has $m$ elements, then their cartesian product $A \times B$ has $n \cdot m$ elements.
\end{remark}

\begin{defbox}
   Let $A$ be a set. We say the \textbf{power set of $A$}, denoted $2^A$ or $\mathscr{P}(A)$, is the set of all subsets of $A$, which includes $A$ and the empty set $\emptyset$. 
\end{defbox}

\begin{example}
   Let $A := \{0, 1, 2\}$, then its power set $\mathscr{P}(A)$ is: 
   \[ \mathscr{P}(A) := \{\emptyset, \{0\}, \{1\}, \{2\}, \{0, 1\}, \{0, 2\}, \{1, 2\}, \{0, 1, 2\}\} \]
\end{example}

\begin{remark}
   For each element in a nonempty set, we can either choose to include or not include it in a subset of $A$. This means that for a set $A$ with $n$ elements, its power set has $2^n$ elements.  
\end{remark}

\subsection{Algebra of Set Operations}

We will show our first proposition (a statement which we will prove) here.

\begin{propbox}
   Let $A$, $B$, and $C$ be arbitrary sets. Then the following properties hold:
   \begin{enumerate}
      \item $A \cap (B \cup C) = (A \cap B) \cup (A \cap C)$
      \item $A \cup (B \cap C) = (A \cup B) \cap (A \cup C)$
      \item $A \setminus (B \cup C) = (A \setminus B) \cap (A \setminus C)$
      \item $A \setminus (B \cap C) = (A \setminus B) \cup (A \setminus C)$
      \item Given an ambient set $U$ such that $A \subseteq U$, $(A^c)^c = A$
   \end{enumerate}
\end{propbox}

\begin{proof}
   Each statement in this proposition asserts that two sets are equal. By Remark 1.0.6, we prove each equality by showing that each side is a subset of the other.

   Let us start by proving (1). We first show that $A \cap (B \cup C) \subseteq (A \cap B) \cup (A \cap C)$. Let $x \in A \cap (B \cup C)$. By definition, $x \in A$ and $x \in B \cup C$. Therefore, $x \in B$ or $x \in C$.

   If $x \in B$, then $x \in A \cap B$, and hence $x \in (A \cap B) \cup (A \cap C)$. If $x \in C$, then $x \in A \cap C$, and hence $x \in (A \cap B) \cup (A \cap C)$. Therefore, \[A \cap (B \cup C) \subseteq (A \cap B) \cup (A \cap C).\]

   Conversely, let $x \in (A \cap B) \cup (A \cap C)$. Then $x \in A \cap B$ or $x \in A \cap C$.

   If $x \in A \cap B$, then $x \in A$ and $x \in B$, so $x \in A \cap (B \cup C)$. If $x \in A \cap C$, then $x \in A$ and $x \in C$, so $x \in A \cap (B \cup C)$. Therefore, \[(A \cap B) \cup (A \cap C) \subseteq A \cap (B \cup C).\] Hence condition (1) holds.

   For (2), let $x \in A \cup (B \cap C)$. If $x \in A$, then $x \in A \cup B$ and $x \in A \cup C$. Therefore, $x \in (A \cup B) \cap (A \cup C)$.

   If $x \in B \cap C$, then $x \in B$ and $x \in C$. Thus $x \in A \cup B$ and $x \in A \cup C$, so again $x \in (A \cup B) \cap (A \cup C)$. Hence \[A \cup (B \cap C) \subseteq (A \cup B) \cap (A \cup C).\]

   Conversely, let $x \in (A \cup B) \cap (A \cup C)$. Then $x \in A \cup B$ and $x \in A \cup C$.

   If $x \in A$, then $x \in A \cup (B \cap C)$. Otherwise, $x \notin A$. Since $x \in A \cup B$ and $x \notin A$, we must have $x \in B$. Similarly, since $x \in A \cup C$ and $x \notin A$, we must have $x \in C$. Thus $x \in B \cap C$, and therefore $x \in A \cup (B \cap C)$. Hence \[(A \cup B) \cap (A \cup C) \subseteq A \cup (B \cap C),\] so condition (2) holds.

   For (3), let $x \in A \setminus (B \cup C)$. By definition, $x \in A$ and $x \notin B \cup C$. Since $x \notin B \cup C$, we have $x \notin B$ and $x \notin C$; otherwise, $x$ would belong to $B \cup C$. Thus $x \in A \setminus B$ and $x \in A \setminus C$. Therefore, \[x \in (A \setminus B) \cap (A \setminus C),\] and hence \[A \setminus (B \cup C) \subseteq (A \setminus B) \cap (A \setminus C).\]

   Conversely, let $x \in (A \setminus B) \cap (A \setminus C)$. Then $x \in A \setminus B$ and $x \in A \setminus C$, which means $x \in A$, $x \notin B$, and $x \notin C$. Since $x \notin B$ and $x \notin C$, it follows that $x \notin B \cup C$. Thus $x \in A \setminus (B \cup C)$. Therefore, \[(A \setminus B) \cap (A \setminus C) \subseteq A \setminus (B \cup C),\] and condition (3) holds.

   For (4), let $x \in A \setminus (B \cap C)$. By definition, $x \in A$ and $x \notin B \cap C$. Therefore, $x \notin B$ or $x \notin C$.

   If $x \notin B$, then $x \in A \setminus B$. If $x \notin C$, then $x \in A \setminus C$. Hence $x \in (A \setminus B) \cup (A \setminus C)$. Therefore, \[A \setminus (B \cap C) \subseteq (A \setminus B) \cup (A \setminus C).\]

   Conversely, let $x \in (A \setminus B) \cup (A \setminus C)$. Then $x \in A \setminus B$ or $x \in A \setminus C$. Thus $x \in A$ and either $x \notin B$ or $x \notin C$. Therefore, $x \notin B \cap C$, so $x \in A \setminus (B \cap C)$. Hence \[(A \setminus B) \cup (A \setminus C) \subseteq A \setminus (B \cap C),\] and condition (4) holds.

   Finally, by definition of the complement relative to $U$, \[A^c = U \setminus A.\] Therefore, \[(A^c)^c = U \setminus (U \setminus A) = A.\] Hence condition (5) holds.
\end{proof}

\begin{propbox}[De Morgan's Laws]
   Let $U$ be a set, and let $B,C \subseteq U$. Then:
   \begin{enumerate}
      \item $(B \cap C)^c = B^c \cup C^c$.
      \item $(B \cup C)^c = B^c \cap C^c$.
   \end{enumerate}
\end{propbox}

\begin{proof}
   For (1), let $x \in (B \cap C)^c$. By definition of the complement, $x \in U$ and $x \notin B \cap C$. Therefore, it is not the case that both $x \in B$ and $x \in C$. Hence $x \notin B$ or $x \notin C$.

   If $x \notin B$, then $x \in B^c$. If $x \notin C$, then $x \in C^c$. In either case, $x \in B^c \cup C^c$. Thus \[(B \cap C)^c \subseteq B^c \cup C^c.\]

   Conversely, let $x \in B^c \cup C^c$. Then $x \in B^c$ or $x \in C^c$, so $x \notin B$ or $x \notin C$. Therefore, $x \notin B \cap C$. Since $x \in U$, it follows that $x \in (B \cap C)^c$. Hence \[B^c \cup C^c \subseteq (B \cap C)^c.\] Therefore, \[(B \cap C)^c = B^c \cup C^c.\]

   For (2), let $x \in (B \cup C)^c$. By definition, $x \in U$ and $x \notin B \cup C$. Therefore, $x \notin B$ and $x \notin C$. Thus $x \in B^c$ and $x \in C^c$, which implies that \[x \in B^c \cap C^c.\] Hence \[(B \cup C)^c \subseteq B^c \cap C^c.\]

   Conversely, let $x \in B^c \cap C^c$. Then $x \in B^c$ and $x \in C^c$, meaning that $x \notin B$ and $x \notin C$. Therefore, $x \notin B \cup C$. Since $x \in U$, we conclude that $x \in (B \cup C)^c$. Thus \[B^c \cap C^c \subseteq (B \cup C)^c.\] Therefore, \[(B \cup C)^c = B^c \cap C^c. \qedhere \]
\end{proof}
